\documentclass[preview]{standalone}

\usepackage{amsmath}
\usepackage{amssymb}
\usepackage{stellar}
\usepackage{definitions}

\begin{document}

\id{psicologia-funzioni-comunicazione}
\genpage

\section{Funzioni della comunicazione}

\begin{snippetdefinition}{numero-dunbar-definition}{Numero di Dunbar}
    Il \emph{numero di Dunbar} indica la grandezza media dei gruppi umani,
    stimata intorno a \(140\)-\(150\) persone.
\end{snippetdefinition}

\begin{snippet}{ultrasocialita-comunicazione}
    La specie umana è ultrasociale: vive in gruppi estesi e dipende dalla
    comunicazione per costruire, mantenere e modificare legami sociali. Per
    questo la comunicazione è centrale nell'esistenza individuale e collettiva.
\end{snippet}

\begin{snippet}{funzioni-comunicazione-overview}
    La comunicazione svolge contemporaneamente tre funzioni principali:
    proposizionale, relazionale ed espressiva.
\end{snippet}

\includesnpt[width=58\%,src=/snippet/static-psychology/communication-functions.jpg]{centered-img}

\begin{snippetdefinition}{funzione-proposizionale-definition}{Funzione proposizionale}
    La \emph{funzione proposizionale} della comunicazione consente di elaborare,
    organizzare, trasmettere e condividere conoscenze facendo ricorso a
    proposizioni.
\end{snippetdefinition}

\begin{snippet}{funzione-proposizionale-note}
    La funzione proposizionale è simbolica e permette elevati livelli di
    astrazione. Grazie a essa è possibile comunicare conoscenze dichiarative,
    collegare informazioni e accrescere il patrimonio delle conoscenze tramite
    apprendimento e memoria.
\end{snippet}

\begin{snippetdefinition}{funzione-relazionale-definition}{Funzione relazionale}
    La \emph{funzione relazionale} della comunicazione consiste nel dare forma
    alla trama delle relazioni interpersonali nei rapporti tra noi e gli altri.
\end{snippetdefinition}

\begin{snippet}{funzione-relazionale-note}
    La comunicazione permette di generare, sviluppare, mantenere, rinnovare,
    cambiare, restaurare o estinguere una relazione. Ogni rapporto sociale viene
    quindi costruito e regolato anche attraverso processi comunicativi.
\end{snippet}

\begin{snippetdefinition}{funzione-espressiva-definition}{Funzione espressiva}
    La \emph{funzione espressiva} della comunicazione è la modalità attraverso
    cui si manifestano pensieri, emozioni, valori e ideali.
\end{snippetdefinition}

\begin{snippet}{funzione-espressiva-note}
    La funzione espressiva è alla base della creatività umana, perché consente
    di rendere visibili e condivisibili stati interni, significati personali e
    forme originali di esperienza.
\end{snippet}

\begin{snippetdefinition}{comunicazione-definition}{Comunicazione}
    La \emph{comunicazione} è uno scambio interattivo osservabile fra due o più
    partecipanti, dotato di un certo grado di consapevolezza e intenzionalità
    reciproca, attraverso cui si condivide un percorso di significati sulla base
    di sistemi convenzionali propri della cultura di riferimento.
\end{snippetdefinition}

\begin{snippetdefinition}{comportamento-definition}{Comportamento}
    Il \emph{comportamento} è qualsiasi azione motoria di un individuo
    osservabile da un altro individuo.
\end{snippetdefinition}

\begin{snippetdefinition}{interazione-definition}{Interazione}
    L'\emph{interazione} è qualsiasi contatto, fisico o virtuale, fra due o più
    individui, anche involontario, capace di modificare lo stato preesistente
    delle cose fra loro.
\end{snippetdefinition}

\begin{snippet}{comunicazione-comportamento-interazione}
    Comunicazione, comportamento e interazione non coincidono. Ogni comunicazione
    è un comportamento, ma non ogni comportamento è comunicazione; ogni
    comunicazione è un'interazione, ma non ogni interazione è comunicazione.
\end{snippet}

\begin{snippetdefinition}{intenzionalita-definition}{Intenzionalità}
    L'\emph{intenzionalità} è sia una proprietà essenziale della coscienza sia
    la proprietà di un'azione compiuta in modo deliberato e volontario per
    raggiungere uno scopo.
\end{snippetdefinition}

\begin{snippet}{intenzione-comunicativa}
    L'intenzione comunicativa consiste nel rendere manifesto al destinatario che
    il parlante vuole comunicargli certi contenuti. Comprende un'intenzione
    antecedente, cioè progettazione e pianificazione, e un'intenzione in azione,
    cioè controllo consapevole durante lo svolgimento dell'attività.
\end{snippet}

\begin{snippetdefinition}{pars-pro-toto-definition}{Pars pro toto}
    La \emph{pars pro toto} è il principio secondo cui, nel manifestare un
    messaggio, il parlante può esprimere solo una parte dei propri contenuti
    mentali.
\end{snippetdefinition}

\begin{snippetdefinition}{opacita-comunicativa-definition}{Opacità comunicativa}
    L'\emph{opacità comunicativa} è la condizione per cui il messaggio espresso
    non rende completamente trasparenti i contenuti mentali e le intenzioni del
    parlante.
\end{snippetdefinition}

\begin{snippetdefinition}{totum-ex-parte-definition}{Totum ex parte}
    Il \emph{totum ex parte} è il principio secondo cui il ricevente attribuisce
    un'intenzione completa al messaggio del parlante sulla base di un numero
    limitato di indizi comunicativi.
\end{snippetdefinition}

\begin{snippetdefinition}{assumere-per-garantito-definition}{Assumere per garantito}
    L'\emph{assumere per garantito} è il principio secondo cui il ricevente tende
    ad accogliere il primo senso del messaggio che gli viene in mente, purché
    non sia immediatamente contraddetto da un altro significato.
\end{snippetdefinition}

\end{document}
